\documentclass[book.tex]{subfiles}

\begin{document}

\chapter{Introduction to Statistical Thermodynamics}
\label{chap:stat_thermo}
\noindent\rule{11cm}{0.4pt} \\
\textbf{Prerequisites:}  \autoref{chap:newtons_laws}, \autoref{chap:electrostatics}, \autoref{chap:classical_thermo}  \\
\textbf{Difficulty Level:} ** \\
\noindent\rule{11cm}{0.4pt}
\vspace{5mm}

Classical mechanics, electromagnetism, and quantum mechanics provide powerful tools to model simple systems with a limited number of microscopic or macroscopic entities. How can we model more complex macroscale systems with a very large number of possibly interacting objects? In this chapter we  introduce the foundations of statistical thermodynamics  which seeks to understand how the macroscopic behavior of systems emerges from the microscopic. 

We refer to a state determined from either quantum or classical mechanics as a \textit{microstate}, to distinguish it from a \textit{state} or \textit{macrostate}, which is a collection of microstates. The macroscopic state of the system is determined by the probability distribution of the microstates, that are accessible to a system through its thermal fluctuations. Our goal then is to find the respective probability of each microscopic state  to the constraints imposed from the perspective of different ensembles. Once the distribution is found, it is possible to determine the statistical information of a macrostate. As we will learn in the rest of this chapter, even when the form of the probability for each microscopic state might change according to its local ensemble, the \textit{expectation value}, which defines a macroscopic observable, remains intact.



\section{Basic Quantities}

An ensemble in statistical mechanics is a probability distribution over a set of microstates such that every such microstate satisfies relations defined in terms of the basic quantities. Figure~\ref{stat_thermo:ensembles} illustrates some ensembles.

\begin{figure}
% \includegraphics[width=\linewidth]{figures/Physics/statistical_mechanics/statistical_thermodynamics/Statistical_Ensembles.png}
\includegraphics[width=\linewidth]{figures/Physics/statistical_mechanics/statistical_thermodynamics/Statistical ensembles.png}
\caption{Illustrations of statistical ensembles. }%TODO: From \url{https://en.wikipedia.org/wiki/Statistical_ensemble_(mathematical_physics)\#/media/File:Statistical_Ensembles.png}. Replace with new figure}
\label{stat_thermo:ensembles}
\end{figure}

The probability distribution of the system is defined over the phase space 

\begin{align}
&p_1,\dotsc,p_n \in \mathbb{R}^3\\
&q_1,\dotsc,q_n \in \mathbb{R}^3
\end{align}
This definition can be directly translated to Physika.

\begin{lstlisting}
ps : $\mathbb{R}^{3}[n]$
qs : $\mathbb{R}^{3}[n]$
ClassicalPhaseSpace[n] = ($\mathbb{R}^3[n]$, $\mathbb{R}^3[n]$)
\end{lstlisting}

It's worth noting that the definitions we've given thus far assume that the phase space of the system has a fixed number of particles $n$. At times, we will want to model open systems that can have a variable number of particles. An alternative definition of the phase space

\begin{lstlisting}
OpenClassicalPhaseSpace[s] = ($n_1$: $\mathbb{N}$,$\dotsc$,$n_s$: $\mathbb{N}$, ClassicalPhaseSpace[$\sum_{i} n_i$] )
\end{lstlisting}

Here there are $s$ kinds of particles and the number of particles of each type is encoded by $n_i$. The type of the phase space is parameterized by the total number of particles 
\begin{align}
n = n_1+\dotsc+n_s
\end{align}
Note that \lstinline{OpenClassicalPhaseSpace} is an example of a dependent type.

Note that \lstinline{ClassicalPhaseSpace} and \lstinline{OpenClassicalPhaseSpace} are very large uncountable spaces. As a result, phase space is often defined into tiny cubes 
\begin{align}
\delta p_1 \times \dotsc \delta p_n \times \delta q_1 \times \dotsc \delta q_n
\end{align}
We can encode this as a type

\begin{lstlisting}
Microstate[$\delta$,$n$] = Interval[$\delta$, $\mathbb{R}$]$^n$
Ensemble[$\delta$, $n$] = ProbabilityDistribution[Microstate[$\delta$, $n$]]
\end{lstlisting}

This discretization process was viewed classically as a mathematical hack, but turns out to have deeper meaning in quantum mechanics. Due to the Heisenberg's uncertainly principle, $\delta$ can be no smaller than the Planck constant $\hbar$.

Statistical mechanical systems are typically parameterized by some macroscopic variables defined on phase space: internal energy $U$, volume $V$, and number of particles $N$, and temperature $T$. We can define some types for these three quantities.

\begin{lstlisting}
System = ClassicalPhaseSpace
U, V, T: System $\to$ $\mathbb{R}$, 
\end{lstlisting}

Statistical mechanics studies a series of relationships between these functions defined on phase space. There are a set of known relationships that govern interactions between common quantities like internal energy, temperature, volume, and number of particles. We say that these values all lie on a common surface defined by
\begin{align}
    f(U, V, N, T) &= 0
\end{align}

The Boltzmann entropy is defined as 
\begin{align}
    S &= k_B \log \Omega
\end{align}
where $\Omega$ is the number of microstates corresponding to a given macrostate. We can parameterize $S$ as a function of other macroscopic variables 
\begin{align}
S(U, V, N, T)
\end{align} 
The joint surface that all these macroscopic quantities lies on satisfies a rich set of relationships.

An important aspect of thermodynamics is that many macroscopic quantities can be expressed in terms of each other. For example, we can reparameterize the internal energy in terms of the entropy of $U(S, V, N, T)$. Many results are derived by shifting the choice of reparameterization.

Another fundamental quantity that we will study is the partition function. 

\begin{lstlisting}
Z: PhaseSpace $\to$ $\mathbb{R}$
\end{lstlisting}

The partition function is arguably the most fundamental quantity in statistical mechanics. As we will learn in future chapters, all other quantities can be derived from the partition function by various formulas.

As a final note, we mention that the physics literature on thermodynamics follows a particular notation. When we write

\begin{align}
\left ( \frac{\partial f}{\partial V} \right )_T
\end{align}
We mean the derivative of function $f$ with respect to $V$ with $T$ held constant. This type of notation is useful for specifying quantities in thermodynamics since we often express quantities in terms of multiple variables and it may not be clear which variables may vary and which must be held constant.
%\textcolor{red}{TODO: What is the type of entropy? Or the partition function?}

%\textcolor{red}{TODO: How can you take the derivative of energy with respect to entropy? The types don't match up?}

\subsection{Extensive and Intensive Variables}

Macroscopic properties in thermodynamics are classified as either extensive or intensive. Extensive quantities depend on $N$, the number of atoms in the system, while intensive quantities don't depend on $N$. Quantities such as mass, volume, and entropy are extensive, while temperature is intensive.

There is a conjugate relationship between extensive and intensive quantities where an extensive change is associated with an associated intensive quantity change. For example, a change in entropy (extensive) is tied to a change in temperature (intensive). 

%\section{TODO}
%\textcolor{red}{TODO: Move these elsewhere}



\end{document}
